\scanpage{27}
\deleted{Structures multiplicatives}

\uline{Exemples} 1) Soient $U$ un anneau, $V$ un sous-anneau, considérons le foncteur
\[
\mathcal C=\mathcal C^U\to\widetilde{\mathcal C}=\mathcal C^V\qquad\text{restriction des scalaires}
\]

C'est un foncteur \uline{exact}, et les applications $\mathrm{Hom}(A,B)\longrightarrow\mathrm{Hom}(\bar A,\bar B)$ sont injectives. De plus, on a
\[
\begin{aligned}
\mathrm{Hom}_{\unsure{V}}(\widetilde A,\widetilde B)&=\mathrm{Hom}_V(A,B)\xrightarrow{\sim}\deleted{\mathrm{Hom}_U(\unsure{A|U\otimes_V A},\unsure{B})}\\
&\inserted{\mathrm{Hom}_U(A,\underbrace{\mathrm{Hom}_V(U,B)})}
\end{aligned}
\]

\margin{$U$-module gauche (grâce à \unsure{$U$}) \deleted{\unsure{la multiplication}} \unsure{via la} \deleted{\unsure{action}} \unsure{multiplication|mult\textsuperscript{on}} de $V$ sur $U$}

cet \deleted{hom.} \inserted{isom} est donné grâce \deleted{au} \unsure{le|à|de} $V$-hom naturel \deleted{(Obtenu en} \deleted{\unsure{faisant} \illegible{1} $B$ \unsure{$\mathrm{Hom}_U(A,B)$|$\mathrm{Hom}_U(U,B)$} $B\to\mathrm{Hom}_V(U,B)$}
\[
e_B:\mathrm{Hom}_V(U,B)\longrightarrow B
\]

défini par \deleted{\unsure{$f$}} $e_B(f)=f(1)$. L'hom. (sur $U$) naturel correspondant
\[
i_B:B\unsure{-|\to}\deleted{\mathrm{Hom}_{\unsure{V}}(U,B)}\mathrm{Hom}_V(U,B)
\]
est celui qui résulte de $B\xrightarrow{\sim}\mathrm{Hom}_U(U,B)\to\mathrm{Hom}_V(U,B)$.

Les suites \deleted{\unsure{$S$}} $\sim$-exactes sont celles qui splittent sur $V$. Les $U$-modules $\sim$-injectifs sont les facteurs directs de modules $\mathrm{Hom}_V(U,B)$.

% paper watermark: ARGENTINE TYPEWRITER (vertical, partly cropped)
% archive sheet number: 26
